Source : Insee, comptes nationaux annuels base 2014 (séries arrêtées) jusqu'en 1995 ; Eurostat et Insee base 2020 ensuite. Calculs de l'auteur (dette/dynamique\_longue.py).\par Lecture : de 1978 à 1995, le ratio de dette passe de 21,2 \% à 56,1 \% du PIB. Les déficits primaires n'y contribuent que pour 8,8 points ; les intérêts (41,9) dépassent nettement la croissance (-33,4), car les taux réels étaient élevés après la désinflation ; les ajustements stock-flux font 17,6 points. De 1995 à 2025, les intérêts et la croissance se compensent et la hausse vient des déficits primaires.
